%% kapa: exam=1 round=37 year=2026 period=1 form=A subject=civil_law subject_ko=민법 seq=11 number=11
%% source: Q-Net 감정평가사 기출문제 (cst00309) artlSeq=5258924 file=raw/1008/5258924/2026년 제37회 감정평가사 1차 1교시.pdf
%% conversion: kordoc 4.24.51 -> tools/gen1.py (auto); review=proofread-sonnet; flags=none
\begin{question}{11}
통정허위표시에 관한 설명으로 옳은 것을 모두 고른 것은? (다툼이 있으면 판례에 따름)
\begin{framedtext}
ㄱ. 민법 제108조 제2항의 통정한 허위의 의사표시의 무효로 대항할 수 없는 제3자는 선의이면 족하고 무과실일 필요는 없다. \par
ㄴ. 가장소비대차의 계약상의 지위를 이전받은 자는 통정허위표시에 따라 형성된 법률관계를 기초로 하여 새로운 법률상 이해관계를 가지게 된 민법 제108조 제2항의 제3자에 해당하지 않는다. \par
ㄷ. 임대차보증금반환채권이 양도된 후 양수인의 채권자가 임대차보증금반환채권에 대하여 압류명령을 받았는데 임대차보증금반환채권 양도계약이 허위표시로서 무효인 경우, 압류채권자는 그로 인해 외형상 형성된 법률관계를 기초로 실질적으로 새로운 법률상 이해관계를 맺은 제3자에 해당한다. \par
\end{framedtext}
\begin{choices}
  \item ㄱ
  \item ㄷ
  \item ㄱ, ㄴ
  \item ㄴ, ㄷ
  \item ㄱ, ㄴ, ㄷ
\end{choices}
\end{question}
\answer{5}
